\scanpage{14}
When $G$ is affine, this should be trivial. We are going to try the

\uline{Proposition} \unsure{Let|Let} $S$ be a quasi-compact prescheme, $n$ an integer $\unsure{\geq1|\geq0}$, \deleted{\unsure{prime|prime}} $S_0=V(n)\subset S$ the closed subscheme defined by $n.1_S$ (thus if $x\in S$, we have $x\in S_0$ iff \deleted{\unsure{characteristic $x$ is of}} \deleted{\unsure{characteristic which divides $n$}} \inserted{characteristic $x\mid n$}), \inserted{$U=S\setminus S_0$,} $G$ an affine group scheme \inserted{of \unsure{finite presentation|finite type}} over $S$, and
\[
(21)\qquad\varphi_U:\mathbb G_a|U\longrightarrow G|U
\]
a group homomorphism. \inserted{We assume $U$ schematically dense in $S$.} Then there exists a \inserted{power} integer \deleted{$d>0$} \inserted{$m=n^{\nu}$ of $n$,} \deleted{such that} and a homomorphism of group schemes
\[
\psi_m:\mathbb G_a\longrightarrow G\unsure{;|,}
\]
such that we have
\[
(22)\qquad\psi_m|U=\bigl(\varphi_U^m:\lambda\longmapsto\varphi_U(\unsure{m\lambda})=\varphi_U(\lambda)^m\bigr)
\]
Let us indeed consider the endomorphism
\[
u_m:\mathbb G_a\longrightarrow\mathbb G_a\qquad\lambda\longmapsto m\lambda
\]
Then formula (22) means that we have commutativity
\diagram{\begin{tikzcd}
\mathbb G_{aU}\arrow[r,dash,"(u_m)_U"]\arrow[dr,"\psi_U"']&\mathbb G_{aU}\arrow[d,"\varphi_U"]\\
&G_U
\end{tikzcd}}
and the existence of the $\unsure{\psi_i|\psi_m}$ for large $m$, means that for large $m$ the composite homomorphism
\[
\deleted{\varphi_U\circ u_m}\qquad\varphi_U\circ u_m:\mathbb G_{aU}\unsure{\longrightarrow|\text{---}}G_U
\]
extends to a group homomorphism $\mathbb G_a\longrightarrow G$.
\margin{13}

% Pied de page imprimé de l'archive.
\begin{center}
\small
13\\
\texttt{https://grothendieck.umontpellier.fr}\\
Université de Montpellier, UMR CNRS 5149, Case courrier 051, Place Eugène Bataillon\\
34095 Montpellier cedex 5 - France
\end{center}

\scanpage{15}
\mspage{655}

For reasons of schematic density and \unsure{uniqueness},
it suffices to prove that $\varphi_U$ extends to a
homomorphism of \deleted{group} schemes -- \inserted{this} extension will
automatically be a homomorphism of group schemes.
\deleted{$A$ \unsure{that of} \illegible{1} $B$}
\inserted{We may assume $S$ \unsure{Affine}, \unsure{let}}
$A$ be \unsure{its} ring, and let $B$ be that of $G$.
So the ring of $U$ will be
\deleted{$A[1/n]\unsure{=}$}
$A_n\overset{\unsure{\mathrm{déf}}}{=}A[1/n]$, that of $G$ will be
$B\otimes_A A_n=B_n$, by hypothesis we have a homomorphism
(corresponding to $\unsure{\varphi_U^*|\varphi_U}$)
\[
(23)\qquad
\diagram{\begin{tikzcd}
B_n \arrow[r,"\varphi"] & A_n[T]
\end{tikzcd}}
\]
whereas the homomorphism \deleted{\unsure{composite}}
\inserted{corresponding} to $\psi_{mU}$ will be the composite of the preceding one
with
\[
(24)\qquad
\diagram{\begin{tikzcd}
A_n[T] \arrow[r,"\unsure{\psi_m^*|u_m^*}"] & A_n[T] \\
T \arrow[r,mapsto] & mT & \text{thus} & T^i \arrow[r,mapsto] & m^iT^i
\end{tikzcd}}
\]
Let \deleted{$\unsure{x}$}
$\inserted{x=}\,a_0+a_1T+\cdots+a_NT^N\in A_n[T]$ \unsure{of which|thus} the
$a_i\in A_n$ we shall therefore have
\[
\psi_m^*(x)=\sum a_i m^iT^i=\sum(a_i n^{\unsure{i\nu}})T^i
\]
thus if $m=n^\nu$ is large enough, we shall have
\[
\deleted{\psi_m^*(x)\unsure{\sim}m^iT^i\in\unsure{A}}
\]
\[
m^i a_i\in A\quad\text{i.e}\quad m^i a_i\in\unsure{A|T^i}[T]
\qquad \forall i\,\uuline{\unsure{\ge|>}1}
\]
thus
\[
\psi_m^*(x)\equiv a_0\pmod{A[T]}
\]
which shows
\[
(25)\qquad \psi_m(x)\in A[T]\ \text{for }m=n^\nu\text{ large}
\quad\Longleftrightarrow\quad a_0\in A
\]
\margin{14}

% Pied de page imprimé de l'archive.
\begin{center}
\small
14\\
\texttt{https://grothendieck.umontpellier.fr}\\
Université de Montpellier, UMR CNRS 5149, Case courrier 051, Place Eugène Bataillon\\
34095 Montpellier cedex 5 - France
\end{center}

\scanpage{16}

Applying this to the images of elements of $B$ which generate $B$ as
an $A$-algebra, we thus have

\uline{Corollary}\quad $S$ \deleted{$\unsure{S_0}$} $n,S_0,U$ being as in
the proposition, and $G$ being an affine scheme of finite type over $S$
(there is no longer any question of a group structure), and
$\varphi_U:\mathbb E^1_U\longrightarrow G_U$ a homomorphism of schemes,
\deleted{for there to \unsure{exist}}
$g_U\,\deleted{\unsure{=}}\overset{\unsure{\mathrm{déf}}}{=}\varphi_U(0)$,
For there to exist a \deleted{\illegible{1}} power $m=n^\nu$
of $n$ such that $\varphi_U\circ\unsure{\psi_m|u_m}$
\deleted{\unsure{(where}} (where $u_m$ given by (24)) extends
\inserted{to $S$}, it is necessary and sufficient that $g_U$ extend to a section $g$ of $G$
over $S$.

\uline{Remark}\quad the preceding arguments prove
in any case the \unsure{tentative} ``theorem'' on page \unsure{662|652}'
in the case where $S$ is noetherian of dimension 1
(the case of interest to us being the one where the generic point of $S$
is of characteristic 0).

\margin{\uline{Group schemes} \uline{associated with}
\unsure{extensions of} \uline{discrete groups}}

3.\quad Let $\unsure{M_0|M}$ now be a free $\mathbb Z$-module
of finite type; we are going to construct a
pro-\uline{group} scheme over $S$, associated
with the projective system of sub-\deleted{groups}\inserted{$\mathbb Z$}-modules
of finite index in $M$. To carry out the
construction, I am going to work out \deleted{\unsure{thus}} a
\margin{15}

% Pied de page imprimé de l'archive.
\begin{center}
\small
15\\
\texttt{https://grothendieck.umontpellier.fr}\\
Université de Montpellier, UMR CNRS 5149, Case courrier 051, Place Eugène Bataillon\\
34095 Montpellier cedex 5 - France
\end{center}

\scanpage{17}
\mspage{656}

\margin{I am \unsure{complicating} my life over the construction of
$\underline{\mathfrak S}$, it suffices to set
$\underline{\mathfrak S}=
\mathfrak S_{S_0}\times_{M_{S_0}}\underline M_{S_0}$.}

preliminary exercise, relative to a situation
\[
(26)\qquad
\diagram{\begin{tikzcd}
1 \arrow[r] & M \arrow[r] & \mathfrak S \arrow[r] & H \arrow[r] & 1
\end{tikzcd}}
\]
exact sequence of discrete groups, with $M$ commutative --
the case which \unsure{will interest} us later on being the one where $M$ is a
free $\mathbb Z$-module of finite type. In any case $M$ is a
$\mathbb Z$-module, and as such defines a Group over
$\deleted{\unsure{\mathbb Z}}\inserted{\unsure{S_0}}=\operatorname{Spec}\mathbb Z$,
namely
\[
\diagram{\begin{tikzcd}
A \arrow[r,mapsto] & M\otimes_{\mathbb Z}A
\end{tikzcd}}
\]
($A$, commutative $\mathbb Z$-algebra) -- which we shall denote by $\underline M$.
In fact, $\underline M$ is representable, i.e. is a group \uline{scheme}
iff $M$ is free of finite type.
I feel like also defining \deleted{schemes in} Groups over $S_0$
(which will be group schemes if $M$ is)
$\underline{\mathfrak S},\underline H$, and an exact sequence
\[
(27)\qquad
\diagram{\begin{tikzcd}
1 \arrow[r] & \underline M \arrow[r] &
\underline{\mathfrak S} \arrow[r] & \underline H \arrow[r] & 1
\end{tikzcd}}
\]
of Groups over $S_0$, and an isomorphism of exact sequences
\[
(28)\qquad
\diagram{\begin{tikzcd}
1 \arrow[r] & M \arrow[r] \arrow[d,"\unsure{\wr}"] &
\mathfrak S \arrow[r] \arrow[d,"\unsure{\wr}"] &
H \arrow[r] \arrow[d,"\unsure{\wr}"] & 1 \\
1 \arrow[r] & \underline M(\mathbb Z) \arrow[r] &
\underline{\mathfrak S}(\mathbb Z) \arrow[r] &
\underline H(\mathbb Z) \arrow[r] & 1
\end{tikzcd}}
\]
we shall take
\[
(29)\qquad \underline H=H_{S_0}
\qquad\text{constant Group with value }H
\]
and the outermost vertical arrows of (28)
\margin{16}

% Pied de page imprimé de l'archive.
\begin{center}
\small
16\\
\texttt{https://grothendieck.umontpellier.fr}\\
Université de Montpellier, UMR CNRS 5149, Case courrier 051, Place Eugène Bataillon\\
34095 Montpellier cedex 5 - France
\end{center}

\scanpage{18}

\unsure{will be|are} the canonical isomorphisms.

We can think of defining $\underline{\mathfrak S}$ as
a functor \inserted{(covariant)} in $A/\mathbb Z$ or (\unsure{contravariant})
in $S/S_0$, by
\[
(30)\qquad
\underline{\mathfrak S}(A)=
\mathfrak S\wedge_M
\underbrace{M(A)}_{\deleted{\unsure{A\otimes_{\mathbb Z}}}\inserted{M\otimes_{\mathbb Z}A}}
\]
\deleted{\unsure{As we have}}.
$\underline{\mathfrak S}(A)$ thus being defined as the extension of
$H$ by $M(A)$ \unsure{deduced} from the extension $\mathfrak S$
of $H$ by $M$ by the homomorphism
\[
(31)\qquad
\diagram{\begin{tikzcd}
M \arrow[r] & M\otimes_{\mathbb Z}A
\end{tikzcd}},
\]
which we consider as a homomorphism of $H$-$\mathbb Z$-modules,
by letting $H$ operate on $M\otimes_{\mathbb Z}A$ via its operation on $M$.

The small difficulty is that the \unsure{group|groups}-functor
thus defined by (30) is not itself representable --
it gives the correct expression for $\underline{\mathfrak S}(A)$ only if
$\operatorname{Spec}A$ is connected -- it corresponds to the subfunctor of
\[
\begin{array}{c}
\inserted{\unsure{\sim}}\\[2pt]
\diagram{\begin{tikzcd}
S \arrow[r,mapsto] &
\unsure{:}\underline{\mathfrak S}(A)=
\deleted{\unsure{\operatorname{Spec}|\underline{\mathfrak S}}}
\underline{\mathfrak S}(S)
\end{tikzcd}}
\qquad (S=\operatorname{Spec}A)
\end{array}
\]
consisting of the \deleted{points such that the}
\inserted{$S\longrightarrow\underline{\mathfrak S}$}
such that the composite
\[
(32)\qquad
\diagram{\begin{tikzcd}
S \arrow[r] & \underline H_S
\end{tikzcd}}
\]
(which corresponds to a \uuline{locally} \uline{\unsure{constant|constant}}
function on $S$ with values in $H$) is
\inserted{(\unsure{entirely})} constant.
\margin{17}

% Pied de page imprimé de l'archive.
\begin{center}
\small
17\\
\texttt{https://grothendieck.umontpellier.fr}\\
Université de Montpellier, UMR CNRS 5149, Case courrier 051, Place Eugène Bataillon\\
34095 Montpellier cedex 5 - France
\end{center}

\scanpage{19}
\mspage{657}

The correct expression will be obtained by starting from the \inserted{\illegible{1}} ``incorrect'' functor
\[
\diagram{\begin{tikzcd}
{[\underline{\mathfrak S}]: A} \arrow[r,mapsto] &
\mathfrak S\wedge_M(M\otimes_{\mathbb Z}A)
\end{tikzcd}}
\]
\unsure{or} better, from the \unsure{contravariant functor|counter-functor}
\[
\diagram{\begin{tikzcd}
S \arrow[r,mapsto] &
\mathfrak S\wedge_M
\Gamma(S,M\otimes_{\mathbb Z}\underline{\unsure{\mathcal O|O}}_S)
\end{tikzcd}},
\]
and taking the associated Zariski sheaf -- in fact we shall have. $=$
\[
(33)\qquad
\underline{\mathfrak S}(S)=
\underset{\mathcal U=(U_i)_{i\in\unsure{S|I}}}{\varinjlim}
\prod_{i\in I}[\mathfrak S](U_i)
\]
\[
\begin{array}{l}
(U_i)_{i\in I}\ \text{a partition}\\
\text{of }S\text{ into }\unsure{\text{schemes}}\\
\unsure{\text{open}}\qquad S\simeq\coprod U_i
\end{array}
\]

Here is another way of presenting the construction of
$\underline{\mathfrak S}$, using the general construction
of §49 -- in the sheaf-theoretic context. We have a system
\[
(34)\qquad
\underbrace{\underline M}_{\text{Group over }S_0},
\quad
\underbrace{\mathfrak S_{S_0}
\overset{\cdot}{\supset}\unsure{M}_{S_0}}_{
\substack{\text{Groups (constant)}\\
\text{over }S_0,\ M_{S_0}\\
\text{normal}}},
\quad
\mathrm{Op.}^{\unsure{\Theta|\theta}}\text{ of }
\mathfrak S_{S_0}\text{ on }\underline M,
\]
\[
\diagram{\begin{tikzcd}
\text{homomorphism }\varphi:
\deleted{\unsure{\mathfrak S}}\inserted{M}_{S_0}
\arrow[r] & \underline M
\end{tikzcd}}
\]
(where the operation of $\mathfrak S_{S_0}$ on $\underline M$,
which \unsure{amounts} to an operation of $\mathfrak S$ on $\underline M$,
is the one deduced from the operation of $\mathfrak S$ on $M$,
by \unsure{functoriality})
these data \unsure{satisfying} the compatibility conditions a) b)
of the passage cited. By the construction \unsure{of the|of} passage cited, we construct
\margin{18}

% Pied de page imprimé de l'archive.
\begin{center}
\small
18\\
\texttt{https://grothendieck.umontpellier.fr}\\
Université de Montpellier, UMR CNRS 5149, Case courrier 051, Place Eugène Bataillon\\
34095 Montpellier cedex 5 - France
\end{center}